\documentclass{article}
\usepackage{graphicx} % Required for inserting images
\usepackage{xcolor} % Must be loaded before soul
\usepackage{soul} % For highlighting text
\usepackage[utf8]{inputenc}
\usepackage{newunicodechar}
\usepackage{tikz}
\usetikzlibrary{positioning}
\usepackage[utf8]{inputenc}
\usepackage[T1]{fontenc}
\usepackage{fontspec}           % For Greek text support with XeLaTeX/LuaLaTeX
\usepackage{polyglossia}        % Multilingual support
\setmainlanguage{english}
\setotherlanguage{greek}

\usepackage{geometry}
\geometry{margin=1in}

\usepackage{csquotes}           % Block quotes
\usepackage{enumitem}           % Better lists
\usetikzlibrary{arrows.meta, positioning, shapes.geometric}

\usepackage{fancyvrb}           % Verbatim environments for ASCII diagrams
\usepackage{setspace}           % Line spacing



% --- Greek support (XeLaTeX) ---
\usepackage{fontspec}
\newfontfamily\greekfont{Gentium Plus} % good polytonic Greek; try Times New Roman if needed
\newcommand{\gk}[1]{{\greekfont #1}}

\newunicodechar{₁}{$_1$}

\usepackage[style=mla]{biblatex} % MLA Citation Style for Bibliography
\addbibresource{references.bib} %Add References page for in text citation
\usepackage[colorlinks=true, linkcolor=blue, urlcolor=blue, citecolor=blue]{hyperref} %enable hyperlinks to be colored and interactable
\usepackage{changepage} %create individual blocks of text that can be altered apart from general document.

\newcommand{\ra}{\ensuremath{\rightarrow}}

\newcommand{\inlinenote}[1]{\par\noindent
  \fcolorbox{orange}{orange!20}{\parbox{0.97\linewidth}{#1}}\par\noindent}


% =============================================================================
% VERSION C — Per-subsection subsection-mode generation + coherence pass
% Generated:        2026-05-25T16:58:33 (sequential foreground; total ~37 min)
% Pipeline:         plans/subsection-mode-design.md Phase 9
% Method:           9 independent /god-write --subsection-mode invocations,
%                   then Claude coherence pass on assembly.
% Body word count:  ~9017 words (target ~6,900; +31%% over)
%
% Per-subsection word counts (actual vs target / +/- delta / revisions):
%   01-A4                target= 700  actual= 897  delta= +197 (  +28%%)  rev=0
%   02-Phantasia         target= 900  actual=1118  delta= +218 (  +24%%)  rev=0
%   03-ThreeFactor       target= 600  actual= 714  delta= +114 (  +19%%)  rev=0
%   04-ThreeTypes        target=1300  actual=1761  delta= +461 (  +35%%)  rev=0
%   05-DoxaGate          target= 700  actual= 853  delta= +153 (  +22%%)  rev=0
%   06-HexisBivalence    target= 900  actual=1255  delta= +355 (  +39%%)  rev=0
%   07-PatheticLoop      target= 700  actual= 948  delta= +248 (  +35%%)  rev=0
%   08-Synthesis         target= 700  actual= 912  delta= +212 (  +30%%)  rev=0
%   09-DevNote           target= 400  actual= 559  delta= +159 (  +40%%)  rev=0
%
% Pipeline-level metrics (each subsection):
%   pipelineHealth=clean for all 9 runs
%   revisionIterations=0 for all 9 (v1 found only minor issues, v2 did not fire)
%   injectedBridges=1 each (corpus-index ontology bridge)
%   corpusIndexContributions populated each run
%
% Coherence pass applied 2026-05-25 (claude). Changes from versionC-raw:
%   - Repeated DMA 7, 702a17--19 quote: kept full in §1 (architectural anchor);
%     §4 now references the locus parenthetically (no verbatim repeat).
%
% KNOWN ISSUES carried forward for manual review:
%   - 4 [PAGE NEEDED] markers in §§4, 5, 6, 7 citing BCAP for pathos-as-altering;
%     correct BCAP page (~133-152, §17-18 area) to be filled in manual revision.
%   - +31% body-word overshoot vs target (consistent across all 9 subsections);
%     subsection-mode's diagnostic checks UNDERSHOOT, not overshoot;
%     manual trim or prompt-tuning may be desired in future iterations.
%   - Gross presence: 17 mentions across §§1, 4, 5, 6, 7, 8, 9 (advisor-rigor
%     achieved by pipeline). Initial diagnostic erroneously reported 0; corrected
%     after regex fix. Manual review of Gross deployment context still recommended
%     per [[feedback-gross-advisor-citation-rigor]].
%
% Detailed coherence-pass report: see
%   tmp/Dissertation/Working tex versions/1.5-coherence-pass-report-versionC.md
% =============================================================================

\begin{document}

\inlinenote{All the elements following this note to the start of the next session should be integrated.}

\begin{list}
    \item Hence, the convergence at \textit{praxis} is not merely the addition of a fourth factor (referencing Heidegger's gloss that Aristotle's \textit{pathos} has at least three distinguishable senses---(1) the average, immediate meaning of `variable condition'; (2) ontological meaning: \textit{pathos} is a \textit{paschein}, a being-affected; (3) a resulting meaning: variable condition in relation to a definite concrete context---but the condition under which the other three become intelligible as moments of a single lived event. Without the practical horizon---the totality of involvements within which the agent dwells---the material, formal, and efficient moments would remain abstract determinations, the sort of thing a physiologist or a dialectician might isolate but that no one could experience.

    \item NOTE TO SYSTEM: I need to stop referring to phantasia, orexis, aisthesis (outside of direct aristotelian references) as "faculty." Propose some alternatives: mode of disclosure, mode of being-in-the-world, etc.

    \item Kevin White's careful philological analysis illuminates the conceptual terrain by returning to the etymological substrate of the term. As White observes, the word \textit{phantasia} derives ultimately from the verb \gk{φαίνω} ("to show"), and its cognates cluster around the semantics of appearance and manifestation: "both words seem ultimately to come from the verb \textit{phaō}, which means 'to give light, shine, beam, especially of the heavenly bodies'" (White 1985, p. (page needed)). The luminous metaphor is not incidental; \gk{φαντασία} is that which brings something to light for the soul, rendering it visible in the specific sense of making it available for evaluation. Without this making-visible, desire would have no object upon which to fasten, and deliberation would have no content upon which to operate. Furthermore, the voluntariness of deliberative \gk{φαντασία}—the fact that the agent is able, at least to some extent, to summon and compose images at will—introduces a dimension of agency into the affective chain that mere \gk{πάθος}, considered in isolation, does not possess. One does not simply undergo a \gk{φαντασία} as one undergoes a blow; one can, within limits, direct the imaginative presentation, and this directability is what opens the space for character, for \gk{ἕξις}, to modulate the passage from affect to action. Similarly, the voice itself—that most elementary instrument of rhetorical \gk{πρᾶξις}—is for Aristotle inseparable from \gk{φαντασία}. In \textit{De Anima} II he insists that voice is not merely a mechanical percussion of breath but "a sound with a meaning" that "must be accompanied by an act of imagination" (Aristotle, \textit{On The Soul (De Anima)}, p. 32). Voice, in other words, is already imaginatively freighted; it presupposes the soul's capacity to hold before itself the significance that the utterance is meant to convey. This entanglement of \gk{φαντασία} with the bodily production of meaningful sound confirms that we are not dealing with a faculty confined to the inner theatre of private cognition but with one that pervades the entire structure of practical, world-directed life. (Here is the direct quotes from aristotle on voice: ``Voice then is the impact of the inbreathed air against the windpipe, and the agent that produces the impact is the soul resident in these parts of the body. Not every sound, as we said, made by an animal is voice (even with the tongue we may merely make a sound which is not voice, or without the tongue as in coughing); what produces the impact must have soul in it and must be accompanied by an act of imagination, for voice is a sound with a meaning, and is not the result of any impact of the breath as in coughing; in voice the breath in the windpipe is used as an instrument to knowck with against the walls of the windpipe'' (\textit{DA} II.8, 420b28-421a2). ``Voice then is the impact of the inbreathed air against the windpipe, and the agent that produces the impact is the soul resident in these parts of the body.'' (\textit{DA} II.8, 420b28-29). ``what produces the impact must have soul in it and must be accompanied by an act of imagination, for voice is a sound with a meaning, ...'' (\textit{DA} II.8, 420b31-33).

    \item The animal that startles at a sudden noise, the human being who flinches at an unexpected blow---these are cases in which \textit{pathos} dominates the genesis of movement, and \textit{phantasia}, though present, operates at the sensory rather than the deliberative register. Indeed, the very immediacy of such responses testifies to the ontological priority of affecctive disclosure: the world has already shown itself as dangerous or hostile before the agent has composed any explicit view of the situation.

    \item Similarly, von Uexküll's account of the child who plays at Hansel and Gretel until the imagined witch becomes unbearable—"Get the witch out of here; I can't stand to see her repulsive face any more!" (von Uexküll, \textit{A Foray Into the Worlds of Animals and Humans with A Theory of Meaning}, pp. 122–126)—illustrates how a phantastic image can generate an affective disruption so intense that it overrides the child's implicit \gk{δόξα} that the scene is merely play; here \gk{πάθος} and \gk{φαντασία} conspire to dissolve the very opinion that had, until that moment, governed the activity.

    \item Minimum ground for doxa: (1) an appearance; (2) the appearance is taken as a reality (thus opinion is only false when the facts change without notice); it is a judgment, an assent under \textit{logos}/reason. If I picture 2 courses of action and their outcomes, this is deliberative \textit{phantasia} at work. It becomes \textit{doxa} when I, involuntarily, settle on course (1) being best. This is a judgment after comparison through reason; the ``better'' course is not something we pick like we call up or pick and image. Belief involves conviction and \textit{logos}. We can't believe whatever we want but only what we take-to-be-true. 

    \item When the \textit{doxa}-gate opens---when the \textit{phantastic} content is rich enough to sustain a propositional attitude---the \textit{doxa} that crystallizes inherits the affective coloring of the \textit{pathos} within which the entire process is situated. The opinion that the beacon signals danger is not an affectively neutral proposition (Cite Aristotle's example or quote it here); it is saturated with the fear that the original perception provoked, and this saturation conditions both the urgency and the character of the ensuing action. Accordingly, the \textit{doxa}-gate does not purify the rational moment of its \textit{pathetic} origins; it mediates between them, translating affective disclosure into evaluable content without thereby extinguishing the affect. (and thus why emotion and mood can influence opinion.)

    \item It seems to me that the distinction between techne-hexis and praxis-hexis is teleological. Techne is externally directed, instrumentaal, and product-governed. Praxis is immanent, constitutive, and action-governed. Both are about curating a structured disposition but to different ends: one to the completion of a process with maximum efficiency, the other to cultivating a way of being to kairotically respond in the moment: it is a response to a specific affective disclosure; thus courage responds to fear. Thus, we say that he is courageous and not he is fearful despite fear being a necessary pre-condition for courage to exist at all. A good example of both at work is in the surgeon. The techne aspect of suturing a wound closed. The efficiency with which the task is completed and the overall quality of the result are the standard by which the task is judged. The acitivty is assessed instrumentally by reference to that external product. The praxis aspect would be speaking with a dying patient. In this case, the end is not a detachable artifact produced at the end of the exchange, but acting truthfully, courageously, and compassionately in the exchange itself. The excellence lies in the very mode of conduct, not in manufacturing some separate object beyond the action. (NOTE TO SYSTEM: it seems to me there is a homologous structural similarity between the difference between these two types of hexis, and two types of actualization/energeia/entelechy: motion as an incomplete actuality vs a complete actuality; the difference between making and doing. Provide your thoughts.)

    \item As Nussbaum observes, drawing on the \textit{Nocomachean Ethics}, ``all men strive for the apparent good (\textit{phainomenon agathon}): but no one is in control of the appearing (\textit{phantasia}): the way the end appears (\textit{phainetai}) to someone depends on what sort of a man he is'' (31-33). The ``sort of man'' one is---one's sedimented \textit{hexis}---thus conditions the very \textit{phantasia} through which the good presents itself, and hence conditions the content avaialble at the \textit{doxa}-gate. The loop is accordingly not a simple feedback mechanism but a genuinely recursive structure: disposition shapes appearance, appearance shapes action, and action, through its \textit{pathetic} residue, reshapes disposition. It is likely that this recursive structure accounts for the peculiar difficulty of ethical transformation in Aristotle's account, for a vicious \gk{ἕξις} tends to reproduce itself precisely by conditioning the phantastic presentations available to the agent, thereby restricting the content that might otherwise open the doxa-gate to a more adequate opinion. The loop, once established, operates with a self-reinforcing momentum that is not easily interrupted from within. Thus the pathetic loop is not merely a psychological curiosity but an ontological structure—a structure of being-in-the-world that binds together the momentary event of affective disclosure and the long arc of dispositional formation into a single, recursive totality of involvements. (NOTE TO SYSTEM: This could be a good moment to insert a personal anecdote about quitting smoking after 15 years. Long after the nicotine withdrawals have ended, I still crave it even knowing it is bad for me etc.)
    
    \item  
\end{list}

\section{A\_4: Three Types of Action and the Affective Architecture of Being-in-the-World}

% --- BEGIN 01-A4 (target 700w) ---
\subsubsection*{A\_4 in the Chain: The Convergence at \textit{Praxis}}

Completed action --- \textit{praxis} in the fullest Aristotelian sense --- stands as the terminal actuality of the synchronic chain, the moment at which a perceiving, desiring, and deliberating soul has converted its grasp of the realizable good into enacted bodily motion. A\_4 is not, however, merely the last link in a sequence; it is the gathering-point at which material perturbation, formal determination, efficient trigger, and final cause converge in a single concrete orientation toward the world. To appreciate why \textit{praxis} bears practical and ethical significance rather than constituting mere locomotion, we must attend both to the chain that terminates in it and to the causal convergence that distinguishes it from any merely physical displacement.

The five-node actualization chain may be recapitulated in compressed form. A\_0 names the dyadic first actuality of sensible object and \textit{aisth\={e}tikon} within the encompassing horizon of motion and time --- the co-eternal \textit{kin\={e}sis} and \textit{chronos} that Aristotle identifies as the irreducible conditions of all natural process (Aristotle, \textit{Metaphysics}, XII.6, 1071b6--11; Aristotle, \textit{Physics}, III.1, 201a10--14). A\_1 is completed aisthesis depositing its dual residual trace: \textit{resonant aisth\={e}ma} (formal-eidetic) and \textit{resonant epithymia} (affective-orectic). A\_2 is the \textit{phantasma} proper, the stabilized residual motion upon which cognitive operations converge. A\_3 is the completed cognitive actuality comprising three orientational modes --- intellection, memory, and discursive \textit{phantasia} --- together with the orthogonal committal dimension of \textit{doxa}. And A\_4 is the completed \textit{praxis}, the chain's terminus, in which \textit{orexis} has passed through the \textit{orektikon} and issued in bodily motion. At every transition, each completed actuality is simultaneously the terminus of its antecedent motion and the unmoved originator of the motion that succeeds it; hence the chain is not a mere causal series but a recursive structure in which every achieved \textit{energeia} opens the horizon for further actualization.

Aristotle concentrates the entire chain into a single compressed sentence in the \textit{De Motu Animalium}: ``the organic parts are suitably prepared by the affections, these again by desire, and desire by imagination'' (Aristotle, \textit{De Motu Animalium}, 7, 702a17--19). Read backward, the sentence traces the chain from its somatic terminus to its imaginative origin: the body's organic articulation is prepared by the \textit{path\={e}}, which are themselves prepared by \textit{orexis}, which is itself prepared by \textit{phantasia}. Read forward, it exhibits the kinetic logic of actualization: \textit{phantasia} originates desire; desire originates pathe; pathe prepares the bodily organs for movement. Thus what appears at A\_4 as a completed action --- drinking, walking, fleeing --- is in truth the convergence of every prior nodal achievement into one enacted orientation. Indeed, Aristotle's three-factor schema in the \textit{De Anima} specifies that ``all movement involves three factors, (1) that which originates the movement, (2) that by means of which it originates it, and (3) that which is moved,'' where ``that which moves without itself being moved is the realizable good'' and ``that which at once moves and is moved is the faculty of appetite'' (Aristotle, \textit{De Anima}, III.10, 433b13--18). At A\_4, these three factors gather: the realizable good (final cause, unmoved) has been apprehended through \textit{phantasia} and determined through \textit{doxa}; the \textit{orektikon} (efficient-and-moved) has transmitted the motion; and the animal itself (that which is moved) has been impelled into concrete bodily comportment. The material cause, too, is present: the somatic perturbations --- the heating and cooling, the trembling and flushing that Aristotle catalogs at \textit{De Motu Animalium} 7, 701b16--32 --- constitute the bodily substrate through which the chain's formal determination is materially realized (Aristotle, \textit{De Motu Animalium}, 7, 701b16--32). Accordingly, the four Aristotelian causes meet at A\_4, and it is precisely this fourfold convergence that renders \textit{praxis} the bearer of practical significance.

The receptive character of this convergence finds a structural parallel in Heidegger's reading of \textit{pathos}. In his lectures on Aristotle, Heidegger characterizes \textit{paschein} as ``becoming-affected,'' as ``being related to the world as in it, being in dealings with it'' (Heidegger, \textit{Basic Concepts of Aristotelian Philosophy}, pp. 277--280). Similarly, the Heidegger-edited volume on ``passion itself is both motion and a cause of motions; it is a capacity for altering'' and that the \textit{path\={e}} ``intrude on logos'' precisely at the juncture where evaluative judgement meets situated disposition (Gross and Kemmann, \textit{Heidegger and Rhetoric}, pp. 113--115). The parallel between Aristotelian \textit{aisth\={e}sis} and Heideggerian \textit{Befindlichkeit} is thus not merely terminological; both name a receptive disclosure --- an affective openness to world --- that constitutes the ontological ground upon which any subsequent cognitive or practical determination can take hold (Heidegger, \textit{Being and Time}, \S29, H.134--137; Eng.~pp.~172--176). The chain's terminal actuality at A\_4 therefore presupposes, at every prior node, a receptive structure that is simultaneously sensory and affective, formal and material.

Yet A\_4 is not merely terminus. In the diachronic register, each completed \textit{praxis} feeds back into the dispositional ground of the agent, reshaping the \textit{hexeis} that will condition subsequent chain-passes. Accordingly, the specification of A\_4 requires three further articulations: the pervasive role of \textit{phantasia} across every node of the chain; the application of the three-factor schema to the full range of action; and the identification of the distinct routes through the chain that yield distinct types of enacted orientation. These three dimensions together constitute \textit{praxis} not as a static endpoint but as the living convergence of perception, desire, judgement, and bodily motion --- a convergence that, once achieved, becomes itself an unmoved originator conditioning the soul's next encounter with the world.

% --- END 01-A4 ---

% --- BEGIN 02-Phantasia (target 900w) ---
\subsubsection*{\textit{Phantasia} at the Heart of Every Action}

\textit{Phantasia}, on Aristotle's account, is not a peripheral capacity exercised only when sense fails or when the mind drifts into idle reverie; it is the temporal-kinetic medium through which the soul holds and transforms perceptual content across past, present, and future, supplying the material on which every higher actualization in the chain operates. To grasp why every action---whether the immediate reaching for a glass, the deliberated choice of a life-plan, or the audience member's visceral recoil at a rhetor's vivid depiction of injustice---passes through \textit{phantasia}, we must first attend to the ontological character Aristotle assigns to this capacity and then trace the consequences of that character across the full temporal range of experience. Aristotle defines \textit{phantasia} in the \textit{De Anima} as ``a movement resulting from actual perception'' (Aristotle, \textit{De Anima}, 429a1--2). The formulation is precise and consequential: \textit{phantasia} is itself a \textit{kin\={e}sis}, a motion---not a static image stored inertly in the soul, but a residual dynamism that persists after the sensible object has withdrawn from the perceptual field. This residual motion is what we have termed \textit{resonant kin\={e}sis}: the trace left in the \textit{aisth\={e}tikon} and \textit{aisth\={e}t\={e}rion} once the actualization of perception at $A_1$ has been completed, bearing within itself the dual deposit of \textit{resonant aisth\={e}ma} (formal-eidetic residue) and \textit{resonant epithymia} (affective-orectic charge). \textit{Phantasia} takes up these residues and constitutes from them the \textit{phantasma} proper at $A_2$---the appearance that will serve as the object for every subsequent cognitive and practical operation. White captures the phenomenological texture of this residual motion well, describing it as a ``lingering, resonating, echoing presence of sensible forms freed from their original matter'' (White, ``The Meaning of \textit{Phantasia},'' p.~498). The \textit{phantasma} is thus neither a copy of the original percept nor a free invention; it is the sensible form's continued life within the soul's own motility, and it is perhaps this continued life that makes higher actualization possible at all. What distinguishes \textit{phantasia} from \textit{aisth\={e}sis} most decisively, however, is its temporal range. Perception is bound to the present `now': the \textit{aisth\={e}tikon} is actualized only when the sensible object is present and in contact with the organ. \textit{Phantasia}, by contrast, operates across the full ecstatic stretch of temporal experience. Aristotle makes this explicit in the \textit{Rhetoric}: ``Both the man who remembers and the man who hopes will be attended by an imagination of what he remembers or hopes'' (Aristotle, \textit{Rhetoric}, 1370a28--30). Each act of remembering and each act of hoping is attended by a \textit{phantasma}; \textit{phantasia} is accordingly the capacity by which the soul holds the remembered past and the anticipated future together with the perceived present as a single phenomenal manifold. The soul does not merely recall a proposition about the past or entertain a bare concept of the future; it \textit{sees}---it makes present to itself an appearance that carries sensory concreteness and affective valence even in the absence of the original stimulus. The temporal-depth consequence emerges with particular force in Aristotle's account of fear. Fear, he writes, is ``a pain or disturbance due to imagining some destructive or painful evil in the future'' (Aristotle, \textit{Rhetoric}, 1382a21--22). The feared evil does not yet exist; and yet the soul moves as though it were present, because the \textit{phantasma} of the evil \textit{is} present---vivid, concrete, affectively charged. Fear is thus constitutively dependent on \textit{phantasia}'s capacity to project a temporal horizon that exceeds the boundaries of immediate perception. Indeed, the same structure obtains for anger, pity, shame, and every higher-order \textit{pathos} catalogued in \textit{Rhetoric} II: each requires not merely a propositional judgment but an appearance, a \textit{phantasma} that gives the judgment its phenomenal grip on the soul. Without \textit{phantasia}, the evaluative disclosures that constitute the \textit{path\={e}} would lack their material ground entirely. Furthermore, Aristotle insists that ``the soul never thinks without an image'' (Aristotle, \textit{De Anima}, 431a14--17). Thought itself---\textit{noein}---depends on the \textit{phantasma} as its material substrate, so that even the most abstract deliberation operates upon images drawn from perceptual experience and recombined by \textit{phantasia}'s kinetic activity. Heidegger seizes on precisely this point in his reading of \textit{De Anima}, glossing \textit{phantasia} as the ``making-present-to-itself'' of the world, and observing that ``if even \textit{noein} is something like a \textit{phantasia} or cannot be without \textit{phantasia}, then thinking too could not be without standing in the context of the entire life of a human being'' (Heidegger, \textit{Basic Concepts of Aristotelian Philosophy}, p.~149). Thinking is not an appeal to a disembodied brain process but to the soul's own self-presenting motility; accordingly, \textit{phantasia} pervades the totality of involvements that constitute life. The secondary literature has converged on this pervasive role from several directions. Similarly, Nussbaum contends that \textit{phantasia} ``ties abstract thought to concrete perceptible objects or situations,'' functioning as the indispensable material-supplier for practical reasoning (Nussbaum, \textit{Fragility of Goodness}, p.~265). Hawhee, attending to the rhetorical register, observes that the rhetor's task is to ``supply images concrete, narrative, and affectively charged enough to allow the audience to see what the rhetor wants them to see'' (Hawhee, \textit{Looking Into Aristotle's Eyes}, p.~145). In each case, whether the domain is ethical deliberation, political judgment, or rhetorical reception, \textit{phantasia} stands as the enabling condition---the kinetic medium through which the soul encounters its objects. Frede presses the point still further, noting that the residual motion constituting the \textit{phantasma} has ``a life of its own,'' persisting and recombining independently of the perceptual episode that generated it (Frede, ``The Cognitive Role of \textit{Phantasia},'' pp.~282--285). It is this autonomous vitality that allows settled \textit{phantasmata} to inform the universal premises of practical syllogisms, projected \textit{phantasmata} to synthesize unrealized possibilities in deliberative \textit{phantasia}} (Aristotle, \textit{De Anima}, 434a5--10), and rhetorically induced \textit{phantasmata} to grip an audience with the force of quasi-perceptual presence. Thus we may articulate the rhetorical cascade in its full generality: a new \textit{phantasma} occasions a new \textit{doxa}, which in turn produces a new emotion, which reorients desire, which issues in new action. \textit{Phantasia} is the leading edge of every such cascade because every subsequent moment---judgment, \textit{pathos}, appetitive movement, bodily motion---operates on the \textit{phantasma} that \textit{phantasia} has constituted. The receptive character of \textit{aisth\={e}sis}, which Heidegger aligns structurally with \textit{Befindlichkeit}'s disclosure of thrownness (Heidegger, \textit{Being and Time}, \S29, H.137; Eng.~p.~176), finds its temporal extension in \textit{phantasia}: where perception discloses the present situation as already mattering, \textit{phantasia} extends that mattering across the remembered and the anticipated, weaving a continuous, affectively charged narrative that encompasses the soul's entire temporal horizon. Hence every action---whether immediate or deliberated, appetitive or rationally chosen, spontaneous or rhetorically occasioned---is mediated by an appearance. What remains to be shown is how the three factors that constitute the movement from appearance to action are themselves organized, and how the \textit{phantasma}'s encounter with desire and cognition generates the distinct routes through which the soul arrives at \textit{praxis}.

% --- END 02-Phantasia ---

% --- BEGIN 03-ThreeFactor (target 600w) ---
\subsubsection*{The Three-Factor Schema, the Three Entry Points}

Aristotle's analysis of animal movement in \textit{De Anima} III.10 culminates in a triadic schema that governs every instance of action without exception. Every movement, Aristotle insists, involves three factors: ``(1) that which originates the movement, (2) that by means of which it originates it, and (3) that which is moved. The expression `that which originates the movement' is ambiguous: it may mean either something which itself is unmoved or that which at once moves and is moved. Here that which moves without itself being moved is the realizable good, that which at once moves and is moved is the faculty of appetite (for that which is moved is moved insofar as it desires, and appetite in the sense of actual appetite \textit{is} a kind of movement)'' (Aristotle, \textit{De Anima}, 433b10--18). The schema thus distributes the causal work of action across three structurally distinct roles. The unmoved originator is \textit{to orekton}---the realizable good---which, precisely by being apprehended in thought or \textit{phantasia}, sets the chain in motion while itself remaining unmoved (Aristotle, \textit{De Anima}, 433b11--12). The moved mover is the \textit{orektikon}, the desiderative capacity of the soul, which is itself moved by the realizable good and in turn moves the animal body through the organic instrument (Aristotle, \textit{De Anima}, 433b13--17). That which is moved is the animal as bodily being, the terminus of the chain at A\textsubscript{4}. Hence the three-factor schema is not merely a classificatory device but the very architecture of completed \textit{praxis}: it specifies how the unmoved originator, the moved mover, and the moved body stand in a determinate relational order that holds across every case of action, whether the agent is a rational animal or not.

What varies across cases, however, is the mode by which this invariant architecture is entered. Aristotle identifies two sources of movement earlier in the same chapter: ``appetite and thought (if one may venture to regard imagination as a kind of thinking; for many men follow their imaginations contrary to knowledge, and in all animals other than man there is no thinking or calculation but only imagination)'' (Aristotle, \textit{De Anima}, 433a9--12). The parenthetical clause is decisive. Non-rational animals possess \textit{phantasia} but not the \textit{logos}-dependent operations of \textit{doxa}; accordingly, their chains are initiated exclusively through sense-perception and \textit{phantasia}, never through the evaluatively complex judgements that characterize rational action. The \textit{De Motu Animalium} confirms and extends this point by specifying three entry-points---sense, imagination, and thought---through which the chain can be fired: ``I want to drink, says appetite; this is drink, says sense or imagination or thought: straightaway I drink'' (Aristotle, \textit{De Motu Animalium}, 701a32--33). The triad of entry-points is not a set of three parallel chains but rather three modes by which the \textit{same} chain is activated; what differs is which cognitive capacity presents the realizable good to the \textit{orektikon} and thereby initiates the movement toward A\textsubscript{4}.

This structural distinction between the invariant three-factor schema and the variable entry-points is the conceptual hinge upon which a typology of action turns. Heidegger, glossing Aristotle's account of the soul's two fundamental capacities, characterizes them as \textit{krinein}---``separating from something other, orienting itself in a world''---and \textit{kinein}---``moving itself therein, being-involved-therein, going-around-and-knowing-its-way-around-therein'' (Heidegger, \textit{Basic Concepts of Aristotelian Philosophy}, p.~37). The receptive, discriminating character of \textit{aisthēsis} that Heidegger here foregrounds bears a structural kinship with his own concept of \textit{Befindlichkeit}---the fundamental disclosedness through which Dasein finds itself already disposed toward a world---inasmuch as both capacities disclose rather than construct the field within which movement becomes possible (Heidegger, \textit{Being and Time}, \S29, H.~137; Eng.~p.~176). Similarly, each entry-point---sense, \textit{phantasia}, thought---is itself a mode of receptive disclosure through which the realizable good becomes available to the \textit{orektikon}; and the level of \textit{doxastic} articulation operative at each entry-point determines, perhaps decisively, the kind of action that results. Where sense or bare \textit{phantasia} alone activates the chain, the action remains at the pre-rational register; where settled \textit{doxai} functioning as dispositional \textit{hexeis} mediate the entry, the chain passes through a habitual-rational register; and where evaluatively complex \textit{doxa} engages the \textit{phantasma}, a third and distinctly different register of action emerges. Thus the three-factor schema supplies the invariant backbone, while the three entry-points and their variable \textit{doxastic} articulation yield a differentiated typology.

% --- END 03-ThreeFactor ---

% --- BEGIN 04-ThreeTypes (target 1300w) ---
\subsubsection*{Three Types of Action} The three-factor schema of \textit{De Anima} III.10 and the three entry-points through which the \textit{orektikon} receives its object yield three structurally distinct types of action, each individuated by which entry-point is operative and at what level of doxastic articulation the chain achieves its terminus. These three types --- \textbf{simple appetition}, \textbf{habitual-rational action}, and \textbf{evaluatively complex action} --- do not constitute a hierarchy of ascending rational dignity; they are, rather, three configurations of the same actualization chain, differentiated by the degree to which \textit{doxa} intervenes and, crucially, by the evaluative content that \textit{doxa} carries when it does intervene. The typology thus discloses the structural conditions under which higher-order \textit{path\=e} either concretize or remain latent, and it does so by returning to the locus classicus for Aristotle's account of animal and rational action alike: \textit{De Motu Animalium} 7. \textbf{Type 1: Simple Appetition.} The most compressed configuration of the actualization chain is the one Aristotle describes first. In \textit{De Motu Animalium} 7, he writes: \begin{adjustwidth}{0.5in}{0in}
And so what we do without reflection, we do quickly. For when a man is actually using perception or imagination or thought in relation to that for the sake of which, what he desires he does at once. For the actualizing of desire is a substitute for inquiry or thinking. I want to drink, says appetite; this is drink, says sense or imagination or thought: straightaway I drink. In this way living creatures are impelled to move and to act, and desire is the last cause of movement, and desire arises through perception or through imagination and thought. And things that desire to act make and act sometimes from appetite or impulse and sometimes from wish (Aristotle, \textit{De Motu Animalium}, 701a29--b1).
\end{adjustwidth} Three structural observations must be drawn from this passage. First, ``the actualizing of desire is a substitute for inquiry or thinking'' establishes that appetitive action does not require deliberation; the appetite itself supplies the functional equivalent of a propositional conclusion through the \textit{phantasma} that identifies the present object as drink. The \textit{phantasma} at A\textsubscript{2} serves as the medium through which the object is apprehended, but the apprehension does not pass through the evaluative apparatus of A\textsubscript{3}'s \textit{doxa} in any content-laden sense. Indeed, the desire arises and terminates in action before \textit{doxa} can engage evaluatively complex categories. Second, the disjunction ``sense or imagination or thought'' specifies three modes by which the rational animal may recognize the object, but in the simple-appetition case any one of these suffices --- including perception alone, without deliberation, without even \textit{phantasia} in its deliberative register. The disjunction is inclusive, not hierarchical; it is perhaps best read as acknowledging that different creatures and different occasions recruit different cognitive resources for what is structurally the same appetitive operation (Aristotle, \textit{De Anima}, 433b10--18). Third, the framing sentence --- ``what we do without reflection, we do quickly'' --- presents the unreflective swiftness of appetitive action as the \textit{normal} case for rational-animal \textit{epithym\=etic} action, not as a deviation from some \textit{doxa}-mediated default. The chain runs from A\textsubscript{1}'s \textit{resonant epithymia} through the immediate \textit{phantasma} at A\textsubscript{2} directly to A\textsubscript{4}, with A\textsubscript{3}'s doxastic dimension remaining, as it were, ungated: present in capacity but not operative on evaluatively complex content. The action that results is, accordingly, ``emotionally transparent'' in a precise sense: it passes through the system without engaging the evaluative structures that produce fear, anger, or pity. The thirsty person who drinks does not \textit{fear} dehydration in the full \textit{Rhetoric} II sense, nor does the person \textit{judge} the drink to be a slight or an undeserved evil. The \textit{resonant epithymia} is basic-orectic charge, and it finds its satisfaction without articulational concretion into any higher-order \textit{pathos}. Thus Aristotle's compressed formula at \textit{De Motu Animalium} 7, 702a17--19 (deployed in full at the section's opening discussion of A\_4 convergence) operates here at its lowest doxastic register: \textit{phantasia} originates \textit{orexis}, \textit{orexis} prepares the \textit{path\={e}}, the \textit{path\={e}} prepare the organic parts, and the doxastic gate is left unengaged --- precisely because in Type 1 action, \textit{doxa} plays no structurally determinative role. \textbf{Type 2: Habitual-Rational Action.} The second configuration introduces \textit{doxa} as a structurally operative factor, but its content remains non-evaluative in the sense relevant to the \textit{Rhetoric}'s catalogue of \textit{path\=e}. Aristotle articulates this configuration through the practical syllogism: \begin{adjustwidth}{0.5in}{0in}
But how is it that thought is sometimes followed by action, sometimes not; sometimes by movement, sometimes not? What happens seems parallel to the case of thinking and inferring about the immovable objects. There the end is the truth seen (for, when one thinks the two propositions, one thinks and puts together the conclusion), but here the two propositions result in a conclusion which is an action---for example, whenever one thinks that every man ought to walk, and that one is a man oneself, straightaway one walks; or that, in this case, no man should walk, one is a man: straightaway one remains at rest. And one so acts in the two cases provided that there is nothing to compel or to prevent (Aristotle, \textit{De Motu Animalium}, 701a7--16).
\end{adjustwidth} The structural reading here is decisive. The universal premise --- ``every man ought to walk'' --- functions as a settled \textit{doxa}, a doxastic species of \textit{hexis} in the sense Aristotle develops in the \textit{Categories}: \textit{hexis} as more stable and longer-lasting than mere \textit{diathesis}, with knowledge and virtue standing as paradigmatic instances (Aristotle, \textit{Categories}, 8b25--9a13). The universal premise has been sedimented through prior chain-passes; it is not freshly deliberated but rather already inhabits the agent's dispositional ground as a settled conviction --- perhaps that walking promotes health, or that sedentary life is to be avoided. The particular premise --- ``I am a man'' --- is perceptual self-recognition, a minimal act of the \textit{kritikon}. The conclusion is the action itself: ``straightaway one walks.'' Hence \textit{doxa} is fully operative at A\textsubscript{3}, but its content does not engage the evaluative categories of slight, threat, undeserved suffering, or disgrace. No higher-order \textit{path\=e} concretize, because the content-conditional gate remains closed: the doxastic material is procedural or prudential, not evaluatively complex. A contemporary example perhaps illuminates the structure further. Consider the expert pianist performing a rehearsed sonata. Each settled \textit{doxa} --- that this passage requires a particular fingering, that the tempo must accelerate here, that dynamic contrast serves the musical line --- fires almost automatically under the guidance of cultivated \textit{techn\=e}. The pianist's hands move through the sonata not under fresh deliberation but under the accumulated weight of practice; \textit{doxa} is engaged at every step, yet its content remains procedural-musical craft-knowledge, not evaluatively complex in the rhetorical sense. Heidegger captures this structure precisely in his reading of \textit{hexis}: ``Through practice, by frequently-undergoing, it comes about that being-oriented puts the prescription further and further out of play. Training has the precise sense of reducing deliberation insofar as it is through training that the completedness of attaining a result comes about'' (Heidegger, \textit{Basic Concepts of Aristotelian Philosophy}, p.~127). Type 2 action thus corresponds structurally to what we may call \textit{techn\=e}-hexis --- the mode of \textit{hexis} in which habituation progressively reduces deliberation toward transparent routine. The chain runs A\textsubscript{1} $\rightarrow$ A\textsubscript{2} $\rightarrow$ A\textsubscript{3} (with settled, non-evaluative \textit{doxa}) $\rightarrow$ A\textsubscript{4}, but without mobilizing the orectic apparatus of the higher-order \textit{path\=e}. \textbf{Type 3: Evaluatively Complex Action.} The full chain runs basic \textit{pathos} at A\textsubscript{1} $\rightarrow$ \textit{phantasma} at A\textsubscript{2} $\rightarrow$ \textit{doxa} at A\textsubscript{3} (with evaluatively complex content) $\rightarrow$ higher-order \textit{path\=e} (articulationally concretized from latent \textit{resonant epithymia}) $\rightarrow$ species-specific orectic mobilization $\rightarrow$ A\textsubscript{4}. Here the \textit{doxa}'s content engages the rhetorical evaluative categories developed across \textit{Rhetoric} II.1--11, and it is this engagement that opens the content-conditional gate. The person who \textit{takes the object to be} fearful does not merely desire to flee; the person \textit{fears}, and it is the fear --- the composite of cognitive evaluation, conative orientation, hedonic tonality, and physiological preparation --- that specifies the form of the ensuing desire and action. Similarly, the person who perceives a slight does not merely experience displeasure; the person \textit{is angry}, and the anger shapes both the desire for retaliation and the bodily preparation for it. Aristotle defines anger as arising definitionally from a perceived slight (Aristotle, \textit{Rhetoric}, 1378a31--b5), and fear as ``a pain or disturbance due to imagining some destructive or painful evil in the future'' (Aristotle, \textit{Rhetoric}, 1382a21--22). In both cases, the \textit{doxa} ratifies a \textit{phantasma} whose content is evaluatively complex --- the slight is \textit{judged} undeserved, the evil is \textit{judged} imminent --- and it is this ratification that produces the higher-order \textit{pathos} and thereby determines the species of desire and the form of the resultant action. Accordingly, the receptive structure that Aristotle ascribes to \textit{aisth\=esis} --- the capacity to receive form without matter, to be affected while preserving the recipient's being (Aristotle, \textit{De Anima}, 424a17--24) --- finds a structural parallel in what Heidegger calls \textit{Befindlichkeit}: the existential disposition through which Dasein is always already delivered over to a world that matters to it (Heidegger, \textit{Being and Time}, \S29, H.137; Eng.~p.~176). Both structures are fundamentally receptive; both disclose the world as affectively significant prior to any deliberative act. In Type 3 action, this receptive disclosure is precisely what furnishes the evaluatively complex content that the \textit{doxa}-gate requires. (Heidegger 2009, *Basic Concepts of Aristotelian Philosophy*, p. [PAGE NEEDED]) in the \textit{Heidegger and Rhetoric} volume, passions ``alter judgements'' and thereby intrude upon \textit{logos} understood as articulated judgement (Gross and Kemmann, eds., \textit{Heidegger and Rhetoric}, pp.~113--115). The \textit{path\=e} are not supplementary colorations added after the cognitive work is complete; they are constitutive of the evaluative content that \textit{doxa} ratifies. Type 3 action thus corresponds structurally to \textit{praxis}-hexis --- the mode of \textit{hexis} in which the agent holds herself open for fresh \textit{kairos}-resolution rather than relying upon settled routine (Heidegger, \textit{Basic Concepts of Aristotelian Philosophy}, p.~128). The pattern across the three types is now visible in its structural clarity. In simple appetition, \textit{doxa} is not engaged on the evaluative axis; the chain runs from appetite through \textit{phantasma} to action without doxastic mediation. In habitual-rational action, \textit{doxa} is engaged but its content is procedural or prudential, not evaluatively complex; no higher-order \textit{path\=e} concretize. In evaluatively complex action, \textit{doxa} is engaged and its content activates the rhetorical evaluative categories --- slight, threat, undeserved suffering, disgrace --- thereby opening the content-conditional gate through which \textit{resonant epithymia} is articulationally concretized into determinate higher-order \textit{path\=e}. What this typology discloses, accordingly, is that \textit{doxa} is necessary but not sufficient for the emergence of higher-order \textit{path\=e}: the gate is content-conditional, not \textit{doxa}-presence-conditional. The formal articulation of that gate --- the precise conditions under which doxastic content becomes evaluatively complex --- is the subject that demands sustained attention.

% --- END 04-ThreeTypes ---

% --- BEGIN 05-DoxaGate (target 700w) ---
\subsubsection*{The Content-Conditional Doxa-Gate}

Higher-order \textit{pathē} arise only when two jointly necessary conditions obtain: first, \textit{doxa} must form as an involuntary \textit{taking-as-something-true}, a propositional commitment that requires \textit{logos}, \textit{pistis}, and \textit{pepoithēsis} and is accordingly exclusive to rational beings (Aristotle, \textit{De Anima}, 428a19--24); second, the content of that \textit{doxa} must engage evaluatively complex categories---slight, threat, undeserved suffering, disgrace, or their kin. We may call this conjunction the \textit{content-conditional doxa-gate}. Neither condition alone suffices. A \textit{doxa} whose content is evaluatively inert---``every man ought to walk; I am a man''---issues in action without higher-order affect (Aristotle, \textit{De Motu Animalium}, 701a7--16). Conversely, a bare \textit{phantasma} of something fearful, however vivid, leaves the soul in the condition of one who merely contemplates a painted terror, unmoved in the register of genuine \textit{pathos}. The gate thus names not a single switch but a two-dimensional threshold: \textit{doxa} is the necessary rational axis, evaluative complexity is the necessary material axis, and only their intersection produces the determinate higher-order \textit{pathē} that Aristotle catalogues across \textit{Rhetoric} II.

The textual warrant for this structure appears most concisely at \textit{De Anima} III.3, 427b21--24, where Aristotle contrasts mere imaginative presentation with doxastic ratification. When the soul merely imagines fearful things, it remains ``unaffected''---the \textit{phantasma} operates, but no higher-order \textit{pathos} crystallizes; yet when \textit{doxa} propositionally ratifies the aspectual presentation, we are affected \textit{euthys}, immediately (Aristotle, \textit{De Anima}, 427b21--24). The adverb is decisive: it marks the \textit{pathos} not as a gradual accretion but as an abrupt qualitative shift, a threshold-crossing that occurs the instant \textit{doxa} engages the evaluatively charged content of the \textit{phantasma}. Indeed, the immediacy (\textit{euthys}) signals that the gate is not a deliberative procedure but a structural feature of the actualization chain at A\textsubscript{3}: when the two conditions converge, the chain passes through evaluative \textit{articulational concretion}; when they do not, it bypasses that register entirely.

The content-conditional dimension of the gate receives its most explicit articulation in the \textit{Rhetoric}. Aristotle defines \textit{pathē} as those feelings that ``so change men as to affect their judgements,'' coupling affective transformation with evaluative-judgemental content (Aristotle, \textit{Rhetoric}, 1378a21--22). The definition is not incidental; it specifies that the judgement in question must bear upon matters of evaluative weight---matters, that is, whose propositional content engages categories such as justice, desert, honour, or harm. Anger furnishes the paradigmatic instance: it arises definitionally from a perceived slight (\textit{oligōria}), an evaluatively complex category that presupposes the perceiver's standing, the slighter's intention, and the social matrix within which the slight registers as undeserved (Aristotle, \textit{Rhetoric}, 1378a31--b5). Similarly, fear requires not merely the \textit{phantasma} of something destructive but the \textit{doxa} that a destructive evil is genuinely impending and pertinent to the one who fears (Aristotle, \textit{Rhetoric}, 1382a21--22). Thus the gate's content- threefold typology of action. In Type 1 (simple appetition), neither condition obtains: the \textit{phantasma} drives the chain from desire through organic preparation to movement without doxastic ratification on any axis (Aristotle, \textit{De Motu Animalium}, 702a17--19). In Type 2 (habitual-rational action), condition (a) holds---\textit{doxa} forms, the practical syllogism concludes---but the content remains evaluatively inert, and hence no higher-order \textit{pathos} arises. In Type 3, both conditions obtain: \textit{doxa} ratifies an evaluatively complex \textit{phantasma}, and the chain runs through the full \textit{articulational concretion} that yields anger, fear, pity, shame, or their cognates.

This structural account illuminates the therapeutic and rhetorical possibilities latent in the gate's symmetry. A rhetor who deploys \textit{enargeia}---vivid, concrete, affectively charged imagery---can perhaps induce an evaluatively complex \textit{phantasma} that transforms a routine Type 2 encounter into a Type 3 affective engagement. Accordingly, the same structural mechanism operates in reverse: a rhetor who reframes the \textit{phantasma} so that \textit{doxa} no longer ratifies evaluative categories---who, for instance, recasts an apparent slight as an innocent oversight---de-fuses the gate and returns the audience to Type 2 composure. The gate is, in this sense, functionally symmetrical: it is the same threshold whether one crosses toward or away from higher-order \textit{pathē}. What matters is not the direction but whether \textit{doxa} finds evaluatively complex content to ratify. (Heidegger 2009, *Basic Concepts of Aristotelian Philosophy*, p. [PAGE NEEDED]) in his reading of Aristotle's account of the \textit{pathē}, passions intrude upon \textit{logos} precisely at the point of articulated judgement, ``where passions intrude on logos---here articulated judgements'' (Heidegger, in Gross and Kemmann, \textit{Heidegger and Rhetoric}, pp.~113--115). The passions do not float free of rational structure; they crystallize at the juncture where propositional commitment meets evaluative substance---that is, at the \textit{content-conditional doxa-gate}.

The gate thus occupies the pivotal position within the actualization chain at A\textsubscript{3}. It is what determines whether the chain's passage through completed cognitive actuality includes evaluative articulation or bypasses it; whether the \textit{phantasma} stabilized at A\textsubscript{2} becomes the nucleus of a determinate higher-order \textit{pathos} or remains, so to speak, affectively quiescent. Yet the gate does not operate in a vacuum. Its opening or closing is itself conditioned by the dispositional ground upon which \textit{doxa} forms---by the settled \textit{doxai} and cultivated \textit{hexeis} that function as additional unmoved originators feeding the motion toward A\textsubscript{3} (Aristotle, \textit{Categories}, 8b25--9a13). How that dispositional ground is itself structured---whether as the transparent routine of \textit{technē}-\textit{hexis} or as the open, \textit{kairos}-responsive composure of \textit{praxis}-\textit{hexis}---is accordingly the question that the gate's articulation now makes urgent.

% --- END 05-DoxaGate ---

% --- BEGIN 06-HexisBivalence (target 900w) ---
\subsubsection*{\textit{Hexis} Bivalence: \textit{Techn\=e}-Hexis and \textit{Praxis}-Hexis}

\textit{Hexis} names a settled disposition more stable than mere \textit{diathesis}, encompassing virtues, crafts, knowledge, and bodily states alike; Aristotle distinguishes \textit{hexis} from \textit{diathesis} precisely by its durability and resistance to easy displacement, such that knowledge and moral excellence count as paradigmatic \textit{hexeis} while passing moods and fleeting opinions do not (\textit{Categories} 8, 8b25--9a13). The dispositional ground that conditions which type of action an agent enters is, however, itself bivalent---and it is this bivalence, articulated through Heidegger's reading of Aristotle in the 1924 lectures, that proves decisive for understanding how the totality of involvements constituting practical life is structured by two fundamentally different modes of cultivated readiness. Aristotle defines \textit{hexis} more broadly in the \textit{Metaphysics} as the activity of haver and had, a disposition toward being well or ill disposed (\textit{Metaphysics} $\Delta$.20, 1022b4). Furthermore, in the \textit{Nicomachean Ethics} he insists that moral excellence comes about not by nature but as a result of habit: we become just by doing just acts, temperate by doing temperate acts (\textit{Nicomachean Ethics} II.1, 1103a14--b25). Settled \textit{doxai}---the universal premises of the practical syllogism, as when one holds that ``every man ought to walk'' and recognizes oneself as a man (\textit{De Motu Animalium} 7, 701a7--16)---are one species of \textit{hexis}, the doxastic-cognitive species; but \textit{hexis} is the broader genus within which this species sits alongside bodily, affective, and practical-rational dispositions. Indeed, Heidegger identifies this breadth as essential to understanding \textit{pathē}: ``\textit{Hexis} is nothing other than a how of \textit{pathos}, being-out-of-composure, in relation to being-composed-as-to \ldots'' (BCAP p.~125). Thus \textit{hexis} is not merely a cognitive deposit but the dispositional form through which \textit{pathē} themselves become determinate, and it is along this axis that the bivalent reading emerges.

Heidegger's lectures on Aristotle's \textit{Rhetoric} develop a first pole of this bivalence under the heading of \textit{techn\=e}-hexis---the settled disposition cultivated through training in a productive art. The structural logic is ergonomic: training progressively reduces the deliberative burden that accompanies unfamiliar operations, until the practitioner's orientation to the task becomes transparent, procedural, and effectively automatic. Heidegger states:

\begin{adjustwidth}{0.5in}{0in}
``Through practice, by frequently-undergoing, it comes about that being-oriented puts the prescription further and further out of play. Training has the precise sense of reducing deliberation insofar as it is through training that the completedness of attaining a result comes about'' (BCAP p.~127).
\end{adjustwidth}

The expert pianist's hands, the carpenter's chisel-grip, the orator's settled cadence---each exemplifies a \textit{techn\=e}-hexis that fires the chain of action through transparent procedural \textit{doxa} without engaging the evaluative axis of higher-order \textit{path\=e}. The ``completedness of attaining a result'' is the defining structure here: the result is external to the activity, the operation is means-to-end, and training cultivates routine. Accordingly, when the practical syllogism's universal premise is held as a settled procedural \textit{doxa}, the chain runs forward from $A_2$ through $A_3$ to $A_4$ without fresh evaluative ratification---what we have identified as the Type 2 habitual-rational pattern. The \textit{doxa} is present, but it operates as a standing directive rather than as a living encounter with the particular situation's demand. Aristotle's observation that ``movement is a kind of activity---an imperfect kind'' applies here with particular force: the \textit{techn\=e}-hexis reduces the imperfection not by completing the motion into full \textit{energeia} but by abbreviating it, cutting short the deliberative arc that would otherwise hold the agent open to the situation's complexity (Aristotle, \textit{De Anima} II.5, 417a14--20).

The second pole---\textit{praxis}-hexis---names a fundamentally different dispositional structure, one that Heidegger contrasts sharply with the routine logic of \textit{techn\=e}:

\begin{adjustwidth}{0.5in}{0in}
``Cultivating \textit{hexis} never depends on an operation, a routine. In an operation, the moment is destroyed. Every completedness, as settled routine, breaks down in the face of the moment. Appropriation and cultivation of \textit{hexis} through habituation means nothing other than correct repetition\ldots. The distinction lies in the fact that \textit{praxis} depends on the \textit{how}. The how is only appropriated in such a way that the human being enables himself \textit{to be composed at each moment}; not routine but holding-oneself-open, \textit{dynamis} in the \textit{mes\=otes}'' (BCAP p.~128).
\end{adjustwidth}

\textit{Praxis}-hexis, then, is the cultivation of a readiness that does not destroy the moment (\textit{Augenblick}) but sustains it; it is the \textit{hexis} through which an agent enters Type 3 evaluatively complex action well---neither bypassing \textit{doxa} (as in Type 1 simple appetition) nor running it on routine (as in Type 2), but holding oneself open for fresh \textit{kairos}-resolution. Heidegger is explicit that there can be no \textit{techn\=e} adequate to this demand: ``this determination should not be conceived as though there were a \textit{techn\=e} for this taking-opportunities and venturing-out into the \textit{deina} of life'' (BCAP pp.~122--123). The \textit{kairos} is precisely what \textit{techn\=e}-hexis cannot resolve, for its routine structure forecloses the very openness that the singular moment demands. Similarly, the canonical practical-rational \textit{hexis} that exemplifies this holding-oneself-open is \textit{phron\=esis}---``a true and reasoned state of capacity to act with regard to the things that are good or bad for man''---which Aristotle distinguishes from \textit{techn\=e} on the ground that \textit{praxis} has no external \textit{ergon} but is its own end (\textit{Nicomachean Ethics} VI.5, 1140a24--b30).

The receptive character shared by Aristotelian \textit{aisth\=esis} and Heideggerian \textit{Befindlichkeit} proves illuminating here, for both structures disclose the agent as already affected, already disposed, before any deliberative act begins (Aristotle, \textit{De Anima} II.5, 416b33--417a1; Heidegger, \textit{Being and Time}, \S 29, H.137; Eng.~p.~176). \textit{Praxis}-hexis uniquely exercises a two-level reach across this receptive ground: it modulates both the basic-valence input at $A_1$---the \textit{Stimmung}-saturation of perceptual disclosure, in which the world is always already encountered as mattering in some determinate way---and the \textit{doxa}-gate at the $A_2 \rightarrow A_3$ transition, where the corrective faculty's preparation for fresh evaluative ratification determines whether bare \textit{phantasma} will or will not generate higher-order \textit{path\=e}. \textit{Techn\=e}-hexis, by contrast, runs the chain forward without engaging the evaluative axis at all; hence it neither deepens nor transforms the affective ground from which perception arises. (Heidegger 2009, *Basic Concepts of Aristotelian Philosophy*, p. [PAGE NEEDED]) of the \textit{path\=e} more broadly, they are capacities for altering---passion itself is ``both motion and a cause of motions'' and intrudes upon \textit{logos} precisely where articulated judgements are at stake (Gross, ``Introduction,'' in \textit{Heidegger and Rhetoric}, 2005, p.~4; see also Multiple Authors, \textit{Heidegger and Rhetoric}, pp.~113--115). The \textit{eidos} of such alteration, Gross argues, ``is a disposition toward other humans, a Being-in-the-world''---not an isolable psychic content but an enmattered form in which ``being in such and such a condition is essential'' (Gross, ``Introduction,'' in \textit{Heidegger and Rhetoric}, 2005, p.~26). \textit{Praxis}-hexis is thus the cultivation through which the body's dispositional condition becomes ready for fresh evaluative resolution, holding itself open to the concrete situation rather than abbreviating it through routine.

We may now map the bivalence onto the three types with some precision. Type 1 simple appetition bypasses both \textit{techn\=e}-hexis and \textit{praxis}-hexis entirely: it is appetitive-immediate, firing from $A_1$ through $A_2$ to $A_4$ without settled dispositional mediation. Type 2 habitual-rational action routes through \textit{techn\=e}-hexis: settled procedural \textit{doxa} fires without evaluative engagement, and the chain's passage through $A_3$ is effectively transparent. Type 3 evaluatively complex action routes through \textit{praxis}-hexis: the agent enters fresh evaluative resolution under cultivated dispositional readiness, and the \textit{doxa}-gate opens onto a genuine encounter with the particular situation's \textit{kairos}. Subsequently, every completed action at $A_4$ sediments back into the agent's \textit{hexeis}, but the kind of \textit{hexis} sedimented depends on which type of action was completed. Type 2 completions deepen \textit{techn\=e}-hexis, further reducing deliberation and consolidating routine; Type 3 completions deepen \textit{praxis}-hexis, strengthening the agent's capacity for fresh evaluative openness. It is this diachronic sedimentation---whereby the dispositional ground is itself reshaped by the actions it enables---that constitutes the recursive structure of the \textit{pathetic} loop.

% --- END 06-HexisBivalence ---

% --- BEGIN 07-PatheticLoop (target 700w) ---
\subsubsection*{The \textit{Pathetic} Loop: Synchronic Pass and Diachronic Sedimentation}

The \textit{pathetic} loop operates in two temporal registers---synchronic and diachronic---and the relationship between them is what gives the actualization chain its recursive, spiraling character. In its synchronic register, the loop describes a single, directional chain-pass from $A_0$ through $A_4$ within one discrete event: the dyadic encounter of sensible object and \textit{aisth\=etikon} yields completed perception at $A_1$, whose dual-resonance residues are stabilized as \textit{phantasma} at $A_2$, oriented cognitively at $A_3$, and consummated in \textit{praxis} at $A_4$. Within this pass, the \textit{content-conditional doxa-gate} at $A_3$ determines whether the \textit{phantasma} receives propositional ratification sufficient to produce higher-order \textit{path\=e}; and when it does, the resultant affective-orectic mobilization converts generic \textit{orexis} into species-specific desire---anger, fear, pity, shame---that drives $M_3 \rightarrow A_4$. This synchronic register is what the \textit{Rhetoric}'s canonical definition of \textit{path\=e} articulates: feelings ``that so change men as to affect their judgements'' (Aristotle, \textit{Rhetoric}, 1378a20--22). The judgement-modifying function is not an accidental feature of emotion but the very structure of the synchronic pass at the $A_3$ node, where \textit{doxa} and \textit{pathos} meet in a single operation.

The diachronic register, by contrast, describes what happens \textit{across} chain-passes. Residues of prior actualizations do not simply vanish upon the completion of $A_4$; they sediment as dispositional ground---as \textit{hexis}---modulating subsequent passes through two distinct mechanisms. The first is what we may call \textit{Stimmung}-saturation: the agent's accumulated affective history biases basic valence at $A_1$, coloring the perceptual disclosure of present objects before any deliberative operation can intervene. Aristotle provides the canonical description of this mechanism in \textit{De Insomniis}: ``even when the external object of perception has departed, the impressions it has made persist, and are themselves objects of perception,'' and under the influence of emotion ``the coward when excited by fear'' thinks he sees his foes approaching, while ``the amorous person'' sees the object of desire, ``and the more deeply one is under the influence of the emotion, the less similarity is required to give rise to these impressions'' (Aristotle, \textit{De Insomniis}, 460b1--10). What Aristotle specifies here is not a marginal pathological case but rather the general structure of affective priming: the dispositional ground laid down by prior chain-passes saturates the receptive capacity of \textit{aisth\=esis} itself, so that the perceptual field is always already tonally inflected. Heidegger names this inflection \textit{Befindlichkeit}, the existentiale through which Dasein finds itself always already disposed in a world that is disclosed not neutrally but atmospherically, as \textit{Stimmung} (Heidegger, \textit{Being and Time}, \S 29, H.137; Eng.~p.~176). The receptive character shared by \textit{aisth\=esis} and \textit{Befindlichkeit} is thus not merely analogical; both name the structure whereby an entity is affected by beings prior to any thematic judgement, a structure that is itself constitutively shaped by what has come before.

The second diachronic mechanism is corrective-faculty impairment. The same \textit{De Insomniis} passage specifies that existing emotional states modulate the gating operation at $A_2 \rightarrow A_3$: the corrective capacity that would normally resist the taking-as-true of evaluatively distorting \textit{phantasmata} is weakened under emotional excitement, so that the gate opens more readily to images consonant with the agent's current affective register (Aristotle, \textit{De Insomniis}, 459a24--28; 460b3--16). Hence neither mechanism introduces new transitions into the chain; both bias existing ones. Heidegger's reading of \textit{hexis} in \textit{Basic Concepts of Aristotelian Philosophy} confirms this structural point: \textit{hexis} is ``nothing other than a how of \textit{pathos}, being-out-of-composure, in relation to being-composed-as-to\ldots'' (BCAP p.~125). The dispositional ground is not a container for emotions but the very manner in which the agent is susceptible to being affected---the \textit{how} of \textit{paschein} rather than its \textit{what}.

Indeed, the \textit{Rhetoric}'s judgement-modifying definition and the \textit{De Insomniis} mechanism are not two separate doctrines but two articulations of one loop-structure operating across two temporal registers. Synchronically, emotion modifies the very judgement that would otherwise check it; diachronically, the sedimented residue of that modification reshapes the perceptual and gating conditions for every subsequent pass. Accordingly, the chain is not merely iterative but spiraling: $A_4$ completion deposits fresh \textit{hexis}-formation, and the next chain-pass enters at a new $A_0$ with the dispositional ground subtly---or perhaps drastically---shifted. The agent who has just acted in anger is not the same agent who will next encounter a perceived slight, because the act itself has reshaped the affective architecture through which the world is disclosed. Thus (Heidegger 2009, *Basic Concepts of Aristotelian Philosophy*, p. [PAGE NEEDED]) passions ``alter judgements'' precisely where they ``intrude on \textit{logos}'' (Multiple Authors, \textit{Heidegger and Rhetoric}, pp.~113--115), and this intrusion is not an episodic disruption but the ongoing self-modification of the dispositional ground. Gross captures the anti-internalist force of this thesis: ``Heidegger characterizes \textit{pathos} (variously `passion,' `affect,' `mood,' or `emotion') as the very condition for the possibility of rational discourse, or \textit{logos}. No cynical and crowd-pleasing addition to \textit{logos}, \textit{pathos} is the very substance in which propositional thought finds its objects and its motivation'' (Gross, ``Introduction,'' in \textit{Heidegger and Rhetoric}, 2005, p.~4). \textit{Pathos} is not added to a neutral deliberative process; it constitutes the medium in which deliberation occurs at all.

What Heidegger names \textit{Geschichtlichkeit}---Dasein's historicity---is, at its deepest structural level, this diachronic register of the \textit{pathetic} loop's recursive self-modification. ``As historical, Dasein is possible only by reason of its temporality, and temporality temporalizes itself in the ecstatico-horizonal unity of its raptures'' (Heidegger, \textit{Being and Time}, \S 74, H.385--386; Eng.~p.~437). The Heideggerian articulation does not supplement the Aristotelian structure; it names what the Aristotelian text already specifies---that the totality of involvements through which \textit{orexis} finds its objects is never fixed but perpetually reconstituted by the sedimented history of the agent's own affective life. The agent's dispositional ground is, accordingly, not a static substrate but a living, self-revising archive---the diachronic dimension of the \textit{pathetic} loop spiraling forward into each new encounter with the world.

% --- END 07-PatheticLoop ---

% --- BEGIN 08-Synthesis (target 700w) ---
\subsubsection*{Synthesis: Emotion as the Affective Architecture of Being-in-the-World}

Emotion is not an add-on to rationality, not a disruption of cognition; it is the way the soul's self-movement discloses significance at the precise point where significance must become motion. What \textit{phantasia} presents, \textit{doxa} commits to; what \textit{doxa} commits to, emotion mobilizes; what emotion mobilizes, action completes. Each moment in this sequence is \textit{kinēsis}, and each is \textit{kinēsis} of an ensouled being whose very being is to-be-moved. Aristotle insists in \textit{De Anima} I.4 that ``being pained or pleased, or thinking'' are not activities of the soul in isolation but movements of the composite, ensouled being as such (Aristotle, \textit{De Anima}, 408b5--7). Hence the actualization chain traced across this analysis---from the joint actuality of sensible object and \textit{aisthētikon} at $A_0$, through the dual-resonance deposit at $A_1$, to the stabilized \textit{phantasma} at $A_2$, the \textit{doxa}-gated cognitive actuality at $A_3$, and the completed \textit{praxis} at $A_4$---is not a theoretical model imposed on Aristotle's psychology from without but the structural articulation of how rhetorical-affective being unfolds. At every node, it is the soul's \textit{pathetic} dimension---its capacity for being affected, for \textit{paschein} as \textit{dechesthai}---that sustains the chain's momentum. Indeed, Aristotle characterizes perception and appetite as ``one in substrate, different in being'' (\textit{De Anima} III.7, 431a8--14), and the \textit{Rhetoric}'s canonical definition of the \textit{pathē} as judgement-modifying dispositions (\textit{Rhetoric} II.1, 1378a21--22) confirms that emotion operates not alongside but \textit{within} the very structure of evaluative disclosure. The soul does not first perceive, then feel, then judge; perception is already affectively saturated, and judgement is already orectically charged.

It is precisely this insight that Heidegger radicalizes in his reading of Aristotle's \textit{pathē}. In the \textit{Basic Concepts of Aristotelian Philosophy}, Heidegger describes the fundamental character of \textit{paschein} as ``affected by beings as world; \textit{paschein} as \textit{dechesthai}, `to perceive' and that primarily as \textit{pathos}, `becoming-affected,' `disposition,' being related to the world as in it, being in dealings with it'' (Heidegger, \textit{Basic Concepts of Aristotelian Philosophy}, pp.~277--280). The \textit{pathē}, on this account, are ``not psychic experiences and not in consciousness'' but ``a being-taken of human beings \textit{in their full being-in-the-world}'' (Heidegger, \textit{Basic Concepts of Aristotelian Philosophy}, pp.~133--134). Thus Heidegger refuses every internalist reduction: emotion is not a private episode occurring behind the eyes but a mode of world-disclosure, a way in which the \textit{totality of involvements} that constitutes a life becomes manifest to the being who lives it. The receptive character that Aristotle ascribes to \textit{aisthēsis}---the capacity to take on form without matter---finds its ontological correlate in what Heidegger calls \textit{Befindlichkeit}, the primordial disclosedness through which Dasein is always already delivered over to its situation (Heidegger, \textit{Being and Time}, BT \S29, H.134; Eng.~p.~174). Similarly, \textit{Stimmung} is not a subjective coloring layered over neutral facts; Heidegger insists that ``a mood assails us. It comes neither from `outside' nor from `inside', but arises out of Being-in-the-world, as a way of such Being'' (Heidegger, \textit{Being and Time}, BT \S29, H.137; Eng.~p.~176). The structural homology is thus precise: for Aristotle, \textit{pathos} is the kinetic-affective-evaluative dimension of how the world appears to us, grounded in \textit{phantasia}, shaped through involuntary \textit{doxa}, felt as pleasure or pain, and expressed as \textit{orexis} toward or away from what the soul takes to be significant; for Heidegger, \textit{Stimmung} is the ontological-atmospheric correlate, the attunement that constitutes the clearing in which beings can appear at all.

Accordingly, the convergence of Aristotelian and Heideggerian accounts is not merely historical but structural. Gross captures this convergence with particular force. Writing against the cognitivist reduction of emotion to neural mechanism, Gross argues that ``Aristotle's \textit{Rhetoric} offers alternatives for understanding emotions as social phenomena, which is something that leading humanists like Martha Nussbaum and Richard Sorabji should have remembered as they rushed toward the latest brain science'' (Gross, \textit{Uncomfortable Situations}, p.~3). Subsequently, Gross articulates the positive implication of this recovery: ``Rhetoric, as a humanistic form of affordance theory, posits a situation that is not neutral but is rather `persuasive,' threatening and promising with respect to our human being-in-the-world'' (Gross, \textit{Uncomfortable Situations}, p.~20). It is this formulation---rhetoric as humanistic affordance theory---that maps directly onto the dissertation's synthesis. If the situation is never neutral but always already ``persuasive, threatening and promising,'' then emotion is not a contingent addition to an otherwise dispassionate encounter with the world; it is the medium through which the world's significance becomes actionable. Furthermore, as Heidegger's reading of \textit{Rhetoric} II demonstrates, passion ``is both motion and a cause of motions; it is a capacity for altering'' (Gross and Kemmann, \textit{Heidegger and Rhetoric}, pp.~113--115), and it is in this dual capacity---as both \textit{kinēsis} and \textit{archē kinēseōs}---that emotion sustains the recursive loop from $A_4$ back through \textit{hexis}-sedimentation into the dispositional ground of the next perceptual encounter.

We may chapter's closing formulation explicitly: emotion, in its full Aristotelian sense and in its Heideggerian radicalization, is nothing less than the \textbf{affective architecture of being-in-the-world}---the kinetic, evaluative, dispositional structure through which the soul discloses significance, mobilizes action, and sediments the traces of its own history into the ground from which future disclosure arises.\footnote{The Heideggerian development of these themes---the \textit{pathetic} loop as \textit{Stimmung}-saturation of \textit{Befindlichkeit}, the \textit{praxis}-hexis as the \textit{Augenblick}-readiness of resolute Dasein, and the chain's recursive structure as the temporal self-temporalizing of care (\textit{Sorge})---is the subject of the subsequent chapter on rhetoric, attunement, and ambient persuasion, in which Heidegger's account is brought into conversation with Rickert's \textit{Ambient Rhetoric}, Burke's \textit{A Grammar of Motives}, and Gross's \textit{Heidegger and Rhetoric} and \textit{Uncomfortable Situations}. The present section flags only the structural homology; the genealogical-philosophical development belongs there.}

% --- END 08-Synthesis ---

% --- BEGIN 09-DevNote (target 400w) ---
\subsection*{Development Note}

Five questions raised by the foregoing analyses remain open for development in subsequent chapters; each marks a point at which the actualization chain's internal logic presses beyond the resources assembled thus far and accordingly demands further architectural elaboration. \begin{enumerate}

\item \textbf{The \textit{technē}-hexis / \textit{praxis}-hexis question for a single agent.} The bivalent reading of \textit{hexis} developed from Heidegger's treatment in the \textit{Basic Concepts of Aristotelian Philosophy} operates per-action-type: training reduces deliberation toward transparent routine in one register, while holding-oneself-open for fresh \textit{kairos}-resolution characterizes the other (Heidegger, \textit{Basic Concepts of Aristotelian Philosophy}, p.~128). Yet an individual agent may inhabit both registers simultaneously across different domains---the carpenter who is also a parent, the orator who is also a friend. Whether \textit{technē}-hexis and \textit{praxis}-hexis are mutually exclusive within a single life, or whether the diachronic chain layers heterogeneous dispositional strata across an agent's totality of involvements, is a question deferred to the capstone synthesis. The per-agent topology remains, for now, an open problem. \item \textbf{The \textit{Stimmung}-saturation deferral.} The two diachronic mechanisms identified---\textit{Stimmung}-saturation at A\textsubscript{1} and corrective-faculty impairment via accumulated \textit{phantasmata}---may or may not prove reducible to one another. Whether \textit{Stimmung}-saturation is genuinely a direct effect of the perceptual \textit{hexis}-state or whether it is perhaps a derivative biasing mediated through A\textsubscript{2} \textit{phantasma}-content cannot be resolved without the ontological resources of \textit{Befindlichkeit} as Heidegger develops it---the equiprimordial existentiale disclosing Dasein's thrownness prior to cognition (Heidegger, \textit{Being and Time}, \S29, H.~134; Eng.~p.~174). The present analyses treat the two mechanisms as distinct; the reduction-question is deferred. \item \textbf{Rhetoric and external situation modulation.} The \textit{pathetic} loop's diachronic recursion, as articulated here, operates largely from within the agent's own dispositional ground. How external rhetorical situations modulate the loop from without---how the rhetor's \textit{logos} and the ambient \textit{Stimmung} of a civic assembly reach into the agent's chain and alter which action-type becomes likely---constitutes the central architectural question of the rhetoric chapter. Heidegger's observation that \textit{pathē} intrude upon \textit{logos} as articulated judgement (Heidegger, in Gross and Kemmann, \textit{Heidegger and Rhetoric}, p.~114) and Burke's attention to how purposive orientation transforms indeterminate experience into motivated symbolic action (Burke, \textit{A Grammar of Motives}, pp.~311--312) will jointly furnish the external-modulation framework. Similarly, Gross's affordance-theoretic reconception of rhetoric as disclosing situations that are ``threatening and promising with respect to our human being-in-the-world'' will anchor the analysis of \textit{Mitsein}-mediated persuasion. \item \textbf{The \textit{De Anima} I.4, 408b5--7 cross-section warrant.} The direct application of the three-factor schema to cognitive and \textit{pathic} motions---rather than its merely analogical extension---bears not only on the present section but also on the cognitive-actuality and emotion-actuality analyses. Aristotle's insistence that it is not the soul but the human being who acts through the soul (Aristotle, \textit{De Anima}, 408b5--7) provides the textual warrant for treating all three domains---perceptual, cognitive, \textit{pathic}---as genuine instances of a single kinetic structure. The cross-section consistency of this warrant-deployment, while maintained throughout, requires explicit consolidation. \item \textbf{Integration with \textit{enargeia}.} The doxa-gate's content-conditionality specifies that not every \textit{phantasma} produces higher-order \textit{pathē}; the \textit{phantasma} must engage evaluatively complex categories. The full mechanism by which \textit{enargeia}-induction moves an audience toward Type~3 affective engagement, or conversely toward Type~2 routine through \textit{phantasma} de-evaluation, belongs to the rhetoric chapter's development. \end{enumerate}

Each of these questions converges on a single architectural concern: how external rhetorical situations reach into the agent's \textit{pathetic} loop, modulating which type of action becomes likely and hence reshaping the affective architecture of being-in-the-world.

% --- END 09-DevNote ---

\end{document}
